% Source: https://artofproblemsolving.com/community/c3h1202030_2015_sprint_round_questions
% Stable identity: aops-topic-1202030-q24
All three curves pass through the origin, which satisfies all three equations and so does not count.

\emph{The two circles.} Subtracting the equations gives $-14x-2y=0$, i.e. $y=-7x$; substituting into the first, $50x^2+50x=0$, so the circles meet at $O$ and $(-1,7)$, which is not on $y=-6x$.

\emph{Line and first circle.} $(x-3)^2+(-6x-4)^2=25$ gives $37x^2+42x=0$: $O$ and $x=-\dfrac{42}{37}$, a point not on the second circle.

\emph{Line and second circle.} $(x+4)^2+(-6x-3)^2=25$ gives $37x^2+44x=0$: $O$ and $x=-\dfrac{44}{37}$, a point not on the first circle.

\begin{asy}
usepackage("amsmath");
size(4.5cm);
defaultpen(fontsize(10pt));
draw(circle((3,4),5)); draw(circle((-4,3),5)); draw((-1.6,9.6)--(1.4,-8.4),red);
dot((0,0)); dot((-1,7)); dot((-42/37,252/37)); dot((-44/37,264/37));
label("$O$",(0,0),SE); label("$(-1,7)$",(-1,7),NE);
label("\mbox{}",(min(currentpicture,user=true)+max(currentpicture,user=true))/2,black); // keeps the text after the figure black when LaTeX typesets the labels inline
\end{asy}

So exactly $\boxed{3}$ pairs satisfy exactly two equations.
